% =============================================================================
% cap6 §6.4 Trabalhos futuros
% =============================================================================
\section{Trabalhos futuros}
\label{sec:conclusao-futuros}

A implementação efetiva do sistema no Campus~2 organiza-se em quatro frentes encadeadas, que retomam, na ordem de execução, as lacunas identificadas na Seção~\ref{sec:conclusao-limitacoes}.

A primeira frente é a \textbf{coleta de campo e a construção de um dataset local}. A prioridade imediata é a instalação dos \emph{kits} de câmeras nos pontos P01--P10 caracterizados no Apêndice~\ref{ap:matriz-pontos}, com rotulagem mínima dos primeiros meses de operação. O acúmulo esperado --- da ordem de dezenas de milhares de quadros por trimestre --- permite construir um dataset local cuja distribuição reflita as condições reais de iluminação, distância focal e composição de fauna do Campus, e que servirá tanto à substituição dos elementos fictícios do Capítulo~\ref{cap:resultados} quanto a recursos para a comunidade brasileira de \emph{camera trapping}, hoje sub-representada em datasets públicos.

A segunda frente é o \textbf{\emph{fine-tuning} dos modelos}. Sobre o dataset local consolidado, o caminho de trabalho mais promissor é o \emph{fine-tuning} curto do \modeloSN{} para o subconjunto efetivo de fauna não-alvo do Campus~2, com destaque para o par crítico \emph{F.\,catus}/felinos silvestres de pequeno porte, e a calibração do limiar $\tau$ da Fase~IV sobre indivíduos efetivamente catalogados pela equipe \emph{Gatosdoc2}. O \emph{fine-tuning} do \modeloMD{} é tratado como recurso de reserva, a ser acionado apenas se o desempenho noturno medido ficar consistentemente abaixo do estimado no Capítulo~\ref{cap:resultados}.

A terceira frente é a \textbf{validação empírica e a iteração arquitetural}. Com dataset local e modelos refinados, o protocolo do Capítulo~\ref{cap:metodologia} é reexecutado, substituindo as faixas teóricas do Capítulo~\ref{cap:resultados} pelos valores medidos e classificando cada fase como validada ou ponto de revisão arquitetural. Os pontos de revisão alimentam uma iteração focada: troca pontual de modelo, mudança de regime de \emph{closed-set} para \emph{open-set} contínuo, ou inclusão de heurísticas de pós-processamento por contexto temporal (sequências de quadros do mesmo indivíduo na mesma sessão).

A quarta frente é a \textbf{integração com o \emph{Gatosdoc2} e os desdobramentos de pesquisa}. A integração consiste em conectar o pipeline ao fluxo de trabalho do \emph{Gatosdoc2}, com painel para a equipe (lista de indivíduos detectados, fila de validação humana da Fase~III, candidatos a novo indivíduo da Fase~IV) e canal de comunicação com a comunidade do Campus (relatórios mensais agregados, sem dados pessoais por força da LGPD). Como desdobramento de pesquisa, três linhas se abrem: estudo longitudinal da dinâmica da colônia com base nos registros do sistema, replicação do protocolo em outros campi da USP --- fornecendo evidência multissítio para o desenho proposto --- e ampliação do escopo do classificador para fauna sinantrópica não-felina relevante para vigilância de zoonoses urbanas no contexto brasileiro.
